\section{Data}\label{sec:data}

\subsection{Point-in-time XBRL panels}
The SEC \texttt{companyfacts} endpoint returns every value a registrant has tagged in XBRL since 2009, with the period it describes, the form, the filing date and the accession number. We build a quarterly panel of 38 fields per company from 72 tags (income statement, cash-flow statement, balance sheet, debt by component, impairments, customer deposits and committed credit facilities) together with the contractual debt-maturity schedule. Four problems make raw XBRL unusable for backtests, and each is handled explicitly.

\paragraph{Tag drift.} Companies change tags over time (for example from \texttt{SalesRevenueNet} to the revenue-from-contracts tag introduced with ASC~606). Each field has an ordered list of fallback tags, merged period by period. Total debt is assembled from its current, non-current, commercial-paper and short-term components; operating income is rebuilt as pre-tax income plus interest when it is not tagged.

\paragraph{Fiscal calendars and year-to-date reporting.} Cash-flow statements in 10-Q filings are cumulative, and the fourth quarter appears only inside the 10-K, so facts that share a fiscal-year start are differenced (H1$-$Q1, 9M$-$H1, FY$-$9M). Retailers on 52/53-week calendars report one quarter of 16 or 17 weeks each year (the first quarter at Kroger, the fourth at Costco and AutoZone). The first version of the pipeline accepted only periods of 75 to 105 days and silently dropped that quarter, so trailing-twelve-month sums for these companies covered three fiscal quarters. A quarter is now any period of 75 to 126 days.\ifdefined\AudGapsOld{} Across the \AudLoaded{} companies that load, gaps longer than 126 days in the quarterly spine fell from \AudGapsOld{} (in \AudFirmsOld{} companies) to \AudGapsNew{} (\AudFirmsNew{}), and the last quarter in each panel matches the company's last periodic filing in \AudLatest{} of cases.\fi

\paragraph{Impairments.} Goodwill, intangible and long-lived-asset impairments are added back to operating income, as credit analysts do, because they are non-cash and would otherwise dominate trailing coverage ratios.

\paragraph{Vintages.} For a cutoff $\tau$, the loader keeps only facts \emph{filed} on or before $\tau$ and, for each period, the latest of those filings, so later restatements are excluded. Every model input, benchmark and calibrated parameter is computed from vintages; realized outcomes are measured on the latest data. A company enters a cutoff only if it has 16 quarters of history and its last known quarter ended at most 200 days before $\tau$, which excludes companies that had stopped filing.

Short-term rates, the term spread and the Baa spread come from FRED; none of these series is revised. Nominal GDP, which is revised, is taken from the vintage published on the cutoff date in the Real-Time Data Set for Macroeconomists \citep{croushore2001}. NBER recession dates since 1947 drive the regime model.

\subsection{A survivorship-free universe}
The first version of the paper used the current S\&P~500, which conditions every backtest on survival to 2026. The universe is now built from EDGAR itself. For each calendar year from 2012 to 2024, the XBRL \texttt{frames} API returns one revenue fact for every company that reported one, including companies that were later acquired, delisted or liquidated. A company enters the universe if it reported at least \$50 million of revenue in some year under any revenue tag the loader uses and is not a financial company or REIT (SIC 6000--6799; managed-care insurers, which GICS classifies as health care, are kept). This gives \UniverseN{} companies. At each cutoff, a company is in the sample if its own filings on that date meet the history requirement and its trailing revenue is at least \$50 million. Industries are the Fama--French twelve \citep{fama1997}.

\subsection{A census of bankruptcies}
Item 1.03 of Form 8-K (``Bankruptcy or Receivership'') must be filed within four business days of a Chapter~11 or Chapter~7 petition. Reading every company's complete filing history from the EDGAR \texttt{submissions} API, \UniverseBk{} companies in the universe filed at least one such 8-K. A firm-year has a bankruptcy event if a petition follows the cutoff within eight quarters; firm-years with a petition in the two preceding years are excluded, because the company is already in or just out of bankruptcy. This replaces the convenience sample of 27 prominent 2020 filers used before.

\subsection{The S\&P~500 as of 20 February 2020}
For comparison with the first version, the 2020 analysis is repeated on the index as it stood on 20 February 2020, taken from the Wikipedia constituents table of that day (which lists CIKs) and checked against the independent membership history of \texttt{github.com/fja05680/sp500}; the two lists agree on all 505 tickers.

\section{The twin}\label{sec:model}

\subsection{Stocks, flows and feedback loops}
The model is quarterly and vectorized over Monte Carlo paths. Its stocks are cash, receivables, inventory, payables, customer deposits, net property, plant and equipment, term debt, revolver drawings and equity; a residual ``other net assets'' stock makes the opening balance sheet balance. The accounting identity assets $-$ liabilities $-$ equity $=0$ holds on every path and in every quarter, which is verified by property-based tests over random scenarios.

\paragraph{Demand.} Deseasonalized log revenue growth follows an AR(1) with Student-$t$ innovations and a two-state Markov recession regime $s_t$ \citep{hamilton1989}:
\begin{equation}
g_t=\mu+\phi\,(g_{t-1}-\mu)+\sigma_g\bigl(1+0.5\,s_t\bigr)\varepsilon^{d}_t+\delta\,s_t,\qquad
\mathrm{Rev}_t=B_{t-1}e^{g_t}\,S_{q(t)}\,(1+\xi_t),
\end{equation}
where $S_{q}$ are seasonal factors, $\xi_t$ is the scenario's revenue shock, the regime's entry and exit probabilities are the NBER base rates, and $\delta=-\beta\,\Delta_{\mathrm{GDP}}$ scales the average nominal-GDP shortfall in recessions by the company's estimated revenue cyclicality $\beta$.

\paragraph{Costs and working capital.} Cash operating costs are $F(1+\pi)^t c_t+(v+m_t+\kappa_t)\,\mathrm{Rev}_t$, with fixed costs $F$, a variable-cost ratio $v$, a mean-reverting margin shock $m_t$ correlated with demand, a scenario cost shock $\kappa_t$ and a restructuring multiplier $c_t$. Receivables and inventory adjust to days-based targets with first-order delays (collections and production lag demand); payables and customer deposits adjust immediately, so a collapse in bookings triggers deposit refunds.

\paragraph{Financing.} Debt matures according to the company's own XBRL maturity schedule. A share of maturities is refinanced according to the company's revealed rollover policy, at a rate equal to the short rate plus a term premium and a rating-dependent spread. The treasury raises new debt when cash falls below a target, draws on its committed revolver below a revealed minimum cash balance, and sweeps idle cash above a cap into debt repayment. A \emph{liquidity breach} occurs when cash is negative after all of these.

\paragraph{Feedback loops.} Four loops govern the response to stress (Figure~\ref{fig:loops}):
\begin{itemize}
\item \textbf{R1, leverage--spread spiral.} Trailing interest coverage maps to a synthetic rating and credit spread \citep{damodaran2026}; refinanced and new debt carry that spread, which raises interest expense and lowers coverage further.
\item \textbf{R2, market-access cliff.} The share of a funding need that can be placed falls from 100\% (coverage at least $1.5\times$) to 60\% and 30\% as coverage deteriorates, and to zero while the scenario closes debt markets. Unfunded maturities drain cash.
\item \textbf{B1, cash conservation.} When cash approaches its floor or coverage falls below $2.5\times$, dividends and buybacks stop and capital expenditure falls to maintenance.
\item \textbf{B2, delayed restructuring.} After two quarters below the target margin, fixed costs fall by 4\% per quarter to a floor of $-15\%$, and recover when margins do.
\end{itemize}
R1 and R2 are reinforcing: they turn a revenue shock into a financing shock. B1 and B2 are balancing: they are management's response. Whether a shock fades or spirals depends on which loops dominate, which is the reason for building the model as a system-dynamics model \citep{forrester1961,sterman2000,morecroft2015} rather than as a set of ratios.

\paragraph{Health score.} Each path is summarized by a 0--100 health score that combines liquidity runway (30\%), the synthetic rating (30\%), net debt/EBITDA (20\%) and Altman's Z$''$ \citep{altman2006} (20\%), capped at 25 after a breach. A company's \emph{health score} under a scenario is the median, across Monte Carlo paths, of each path's average score over the horizon, and its \emph{breach probability} is the share of paths with a liquidity breach by the end of the horizon. The \emph{stress health score} uses the demand-freeze scenario and an eight-quarter horizon. The product version of the score also scores coverage and leverage against sector norms; because those norms are estimated from current data, every result in this paper uses the unadjusted score unless stated otherwise.

\subsection{Calibration}
All parameters are estimated from the company's own point-in-time history (Appendix Table~\ref{tab:params}). Seasonal factors are median ratios to a centred moving average; trend growth is the median of the last 12 quarters shrunk halfway toward 1.2\% per quarter; volatilities are robust (median absolute deviation) AR(1) residual standard deviations, with Student-$t$ tails fitted by maximum likelihood. The fixed/variable split comes from a penalized regression of cash operating costs on revenue over 16 quarters, shrunk toward a 35\% fixed share when revenue has not varied enough to identify it. The rollover share is inferred from the two-year debt trend, and the minimum operating cash balance from revealed behaviour (half the 10th-percentile cash balance over three years). A Vasicek short rate \citep{vasicek1977} is fitted to the 3-month Treasury bill. Parameters that are assumptions rather than estimates (for example a credit facility equal to 8\% of revenue when none is tagged) are labelled as such in the software's output.

\subsection{Scenarios}
Eight presets (rates $+200$\,bp, sales $-20\%$, recession, demand freeze, stagflation, credit crunch, supply-chain shock and a baseline) are expanded into per-quarter shock paths. The \emph{demand freeze}, used throughout the evaluation, is a generic 2020-style shutdown: revenue $-70\%$ at the trough for three quarters with an eight-quarter recovery, credit spreads $+400$\,bp, debt markets closed for two quarters and policy rates cut by 150\,bp. The same scenario is applied to every company, so a low stress score measures structural fragility to a common shock, not a forecast of exposure. Reverse stress testing finds the smallest shock that produces a 50\% breach probability, and Sobol indices \citep{sobol2001,saltelli2010} decompose the variance of the health score over the joint scenario space.
